\section{Threats to Validity}\label{sec:threats}
\textbf{Internal.} V7-B was tuned and debugged on the evaluated pairs; the baselines B4 and B5 were fitted on development material; synthetic labels are generator labels, not adjudicated, and at least two are wrong under V7-B's boundary; the BugsInPy units have no truth label; the S1/S2 thresholds are unvalidated defaults; the sufficiency level is a coverage proxy with a known weak link to fault detection~\cite{inozemtseva2014coverage,just2014mut}; one constant-true channel check and the ignorance of input count weaken S2. Label noise and duplication are known to distort evaluations~\cite{allamanis2019,jimenez2019}.
\textbf{Construct.} ``False assurance'' is measured against labels that may not be true divergences; ``abstention'' is dominated by unimplemented features; CPU is child-only.
\textbf{External.} Python top-level functions with JSON arguments; one machine; macOS \texttt{sandbox-exec}; CPython 3.13 only; QuixBugs counts as a single cluster; synthetic projects are generator programs; the effective number of independent projects is not assessable. Leakage and non-independence are expected~\cite{kapoor2023leakage}.
\textbf{Statistical.} Percentile bootstrap intervals undercover with a small effective number of clusters and one dominant cluster; the estimated intraclass correlation of zero is unreliable; power is low (Table~\ref{tab:power}); randomisation-related guidance for evaluations~\cite{arcuri2012a,klees2018evaluating} is only partly applicable. Comparisons are unadjusted exploratory; no multiplicity claim is made.
\textbf{Review.} The literature review is single-reviewer, abstract-level for most sources; absence claims are not proofs.
